%% ====================================================================
%%  It's Good to Be Humble:
%%  Tamarin Prover's Soundness Versus One Evening with an LLM
%%
%%  Build: tectonic -X compile humble.tex
%% ====================================================================
\newif\ifanonymous
\anonymousfalse

\documentclass[version=preprint]{iacrcc}

%% iacrcc already loads amsmath, amssymb, amsthm, mathtools, microtype,
%% hyperref, tikz and graphicx.
\usepackage{fontspec}
\usepackage{enumitem}
\usepackage{booktabs}
\usepackage{tabularx}
\usepackage{float}
\usepackage{placeins}
\usepackage{xcolor}
\usepackage{listings}
\usetikzlibrary{positioning,calc,fit,arrows.meta,backgrounds}
\PassOptionsToPackage{hidelinks}{hyperref}

% Workaround: iacrcc.cls sets the bibliography's PDF anchor with \pdfdest, a
% pdfTeX primitive undefined under XeTeX (which Tectonic uses).
\ifdefined\pdfdest\else
  \def\pdfdest name#1XYZ{}%
\fi

%% --- Fonts --------------------------------------------------------------
%% Fall back gracefully if Berkeley Mono is not installed on the build host.
\IfFontExistsTF{Berkeley Mono}{%
  \setmonofont{Berkeley Mono}[Scale=MatchLowercase]%
}{%
  \setmonofont{DejaVu Sans Mono}[Scale=MatchLowercase]%
}

%% --- Loose-line tolerance ------------------------------------------------
\setlength{\emergencystretch}{3em}
\tolerance=2000
\hbadness=10000
\vbadness=10000

%% --- Float packing --------------------------------------------------------
\renewcommand{\topfraction}{0.92}
\renewcommand{\bottomfraction}{0.65}
\renewcommand{\textfraction}{0.05}
\renewcommand{\floatpagefraction}{0.88}
\setcounter{topnumber}{3}
\setcounter{totalnumber}{5}

%% --- Palette --------------------------------------------------------------
\definecolor{tink}{HTML}{1A1A1A}   % ink
\definecolor{tsoft}{HTML}{6B6B6B}  % secondary text, comments
\definecolor{tline}{HTML}{9A9A9A}  % rules
\definecolor{thot}{HTML}{A35C3A}   % unsound, wrong verdict
\definecolor{tgood}{HTML}{2F6B4F}  % sound, correct verdict
\definecolor{tkey}{HTML}{2B4C7E}   % language keywords
\definecolor{tprim}{HTML}{1F6B70}  % primitives, facts
\definecolor{twash}{HTML}{F4F2EE}  % listing background

%% --- Listings -------------------------------------------------------------
\lstdefinestyle{tbase}{
  basicstyle=\ttfamily\footnotesize\color{tink},
  keywordstyle=\color{tkey}\bfseries,
  keywordstyle=[2]\color{tprim},
  keywordstyle=[3]\color{thot}\bfseries,
  commentstyle=\color{tsoft}\itshape,
  stringstyle=\color{tgood},
  backgroundcolor=\color{twash},
  frame=leftline,
  framerule=0.8pt,
  rulecolor=\color{tline},
  framexleftmargin=4pt,
  xleftmargin=6pt,
  breaklines=true,
  breakatwhitespace=true,
  columns=fullflexible,
  keepspaces=true,
  showstringspaces=false,
  upquote=true,
  aboveskip=6pt,
  belowskip=4pt,
}

\lstdefinelanguage{TamarinL}{
  morekeywords={theory,begin,end,rule,lemma,restriction,builtins,functions,
    equations,All,Ex,let,in,exists-trace,all-traces,not,macros,predicates},
  morekeywords=[2]{Fr,In,Out,K,KU,senc,sdec,aenc,adec,h,pk,last},
  morekeywords=[3]{reuse,sources,use_induction},
  sensitive=true,
  morecomment=[l]{//},
  morecomment=[s]{/*}{*/},
}

\lstdefinelanguage{HaskellL}{
  morekeywords={where,let,in,do,case,of,if,then,else,data,type,instance,
    guard,return,concat,null},
  sensitive=true,
  morecomment=[l]{--},
  morecomment=[s]{\{-}{-\}},
}

\lstnewenvironment{tmlisting}[1][]{\lstset{style=tbase,language=TamarinL,#1}}{}
\lstnewenvironment{hslisting}[1][]{\lstset{style=tbase,language=HaskellL,#1}}{}
\lstnewenvironment{outlisting}[1][]{\lstset{style=tbase,language={},#1}}{}

\newcommand{\code}[1]{\texttt{\small #1}}
\newcommand{\tmc}[1]{\lstinline[style=tbase,language=TamarinL,
  backgroundcolor=\color{white},frame=none,basicstyle=\ttfamily\small]|#1|}

\setlist{topsep=3pt,parsep=0pt,itemsep=2pt}

%% --- Theorem-like ---------------------------------------------------------
\theoremstyle{definition}
\newtheorem{observation}{Observation}
\theoremstyle{remark}
\newtheorem{finding}{Finding}

\title[running   = {It's Good to Be Humble},
       plaintext = {It's Good to Be Humble: Tamarin Prover's Soundness Versus One Evening with an LLM}
      ]{It's Good to Be Humble:\\Tamarin Prover's Soundness Versus\\One Evening with an LLM}

\ifanonymous
  \addauthor[inst={1}]{Anonymous Author}
  \addaffiliation{Anonymous Affiliation}
\else
  \addauthor[inst    = {1},
             email   = {nadim@symbolic.software},
             surname = {Kobeissi}
            ]{Nadim Kobeissi}
  \addaffiliation{Symbolic Software}
\fi

\begin{document}
\maketitle

\keywords{formal verification, protocol analysis, soundness, Tamarin, tool
assurance, LLM-assisted auditing}

\begin{abstract}
The Tamarin prover is among the most theoretically developed tools in symbolic protocol analysis, and it has been used to certify TLS 1.3, 5G-AKA, EMV and Signal. We report five defects in Tamarin 1.12.0 and on the current \code{develop} branch that cause it to return wrong verdicts, all found in a single evening of LLM-assisted auditing by one person who is not a Tamarin developer.

Three defect families survive scrutiny. First, the predicate that decides whether a restriction is a safety formula accepts formulas containing the \code{last()} atom, which are not prefix-closed. This makes Tamarin's induction rule unsound: we exhibit one lemma that Tamarin reports \emph{verified} with induction and \emph{falsified} without, and, by chaining through a \code{[sources]} annotation, a protocol that publishes its key with \tmc{Out(k)} whose secrecy Tamarin reports as verified under a clean wellformedness report. The predicate is byte-identical to the 2012 commit that introduced it to fix bug \#108, and the \code{last} atom was already surface syntax in that same commit. The hole and the fix shipped together, and the hole has been open for fourteen years. Second, the wellformedness checks that catch a lemma quantifying over an action fact no rule emits sit behind a guard that inverts their own docstring, and behind a flag hardcoded to \code{False}; they have not executed since a ProVerif export refactor in PR \#725. Third, lemmas annotated \code{[reuse]} and \code{[sources]} are admitted as assumptions with no check that they were ever proved, and the assumption constrains both trace quantifiers, so one theory makes Tamarin report \emph{found trace} and \emph{no trace found} for the same statement in one summary block.

We report with equal care what we could not break. Tamarin's constraint solver, its term algebra over exclusive-or, Diffie-Hellman, multisets and natural numbers, its newest monotonicity inference, and its stored-proof replay all resisted every probe we aimed at them. The defects are not in the hard part. Every one sits in the trust boundary around the solver, where a check is skipped rather than an inference drawn wrongly. That is the uncomfortable result: the sophistication is real and the guardrails around it are thin, and an evening was enough to find that out.
\end{abstract}

\section{Introduction}\label{sec:intro}

Tamarin~\cite{TamarinCAV,TamarinCSF} is a symbolic protocol verifier with
unbounded sessions, user-definable equational theories, support for stateful
protocols, and a constraint-solving procedure with published soundness and
completeness arguments. It is the tool of record for several of the most
consequential protocol analyses of the last decade: TLS 1.3~\cite{TLS13},
5G-AKA~\cite{FiveG}, EMV~\cite{EMV}, and others. When a standards body wants a
machine-checked argument that a protocol meets its goals, Tamarin is one of two
or three tools, alongside ProVerif~\cite{ProVerif}, that can supply one.

This paper reports five defects that make Tamarin return wrong answers. All
were found in a single evening by one person, working with a
general-purpose large language model as an auditing assistant, on a machine
with no Tamarin build toolchain installed. We did not read the Tamarin papers
in preparation. We did not have prior familiarity with the constraint-solving
implementation. What we had was the source tree, a stock binary from Homebrew,
and about six hours.

We want to be precise about what that does and does not show, because the
temptation to overclaim here is considerable and we have tried to resist it in
both directions.

It does not show that Tamarin is a bad tool, or that its theory is wrong, or
that published Tamarin results are unreliable. Section~\ref{sec:held} reports
at length on the parts we attacked and failed to break, and that list is longer
than the list of defects. Nor does the \code{last()} unsoundness invalidate any
published case study: we scanned all 1042 models shipped with Tamarin and none
of them uses the construct in the position that triggers the bug
(\S\ref{sec:scope}).

What it does show is an asymmetry. Tamarin has been developed for over a
decade by a group with deep expertise, and the machinery that group built held
up under everything we threw at it. The defects we found are all in the layer
around that machinery: whether a formula was classified correctly before being
handed to the solver, whether a wellformedness check runs at all, whether a
lemma admitted as an assumption was ever proved. Those checks are one-line
predicates and list comprehensions. None of them requires a PhD to audit. One
of them has been wrong since August 2012.

\subsection{Why the cost matters}

Finding a soundness bug in a verification tool has historically been expensive.
It has usually meant a researcher with the relevant background, motivated by a
specific suspicion, spending days or weeks. That cost is why tool assurance has
been treated as something that happens occasionally, on purpose, by specialists.

The cost has changed. The audit reported here consisted of reading source with
an assistant that could hold the whole call graph in view, forming hypotheses
about which predicates were load-bearing, and testing each hypothesis with a
ten-line model file. The assistant did not find the bugs by magic and it
proposed several theories that turned out to be wrong (\S\ref{sec:method}
describes two). What it changed was the iteration rate. Reading
\code{isSafetyFormula}, noticing that \code{Last} atoms pass a test named for
existentials, and writing the model that exploits it took under twenty minutes.

If one evening at this cost yields five defects in the most theoretically
developed tool in the field, the reasonable inference is not that this
particular tool is unusually broken. It is that the assumption underneath the
whole practice, that tools of this pedigree have been audited to a depth
proportional to the trust placed in them, was never tested at a cost anyone
could afford. Now it can be.

\subsection{Contributions}

\begin{itemize}
  \item Three confirmed defect families in Tamarin 1.12.0 and on
    \code{develop} at \code{ef3f0468}, each reduced to a self-contained model
    file of at most fifteen lines, each accompanied by a control that isolates
    its mechanism (\S\ref{sec:last}, \S\ref{sec:vacuity}, \S\ref{sec:reuse}).
  \item A fourteen-year provenance analysis of the first defect, establishing
    by direct comparison of commit contents that the flawed predicate is
    byte-identical to the one introduced in 2012 to fix the same class of bug,
    and that the syntax which defeats it was already available in that commit
    (\S\ref{sec:provenance}).
  \item A ground-truth-free auditing method (\S\ref{sec:oracles}) in which a
    tool is caught contradicting itself rather than contradicting a reference
    answer, which is what makes the third defect checkable by a reader who
    knows nothing about the protocol under analysis.
  \item A deliberate account of what resisted attack (\S\ref{sec:held}),
    including the constraint solver, the term algebra, and stored-proof replay,
    which is what locates the defects precisely rather than smearing suspicion
    across the tool.
  \item An artifact: seven model files, a runner that exits non-zero if any
    defect stops reproducing, and two analysis scripts (\S\ref{sec:artifact}).
\end{itemize}

\section{What a Tamarin user is entitled to assume}\label{sec:background}

Tamarin analyzes protocols specified as multiset rewriting rules over a term
algebra, in the Dolev-Yao model~\cite{DolevYao}. A rule consumes facts, emits
action facts onto a trace, and produces facts. Security properties are
first-order formulas over that trace, written as lemmas. The prover answers
\code{verified} or \code{falsified} for each lemma, or reports that its analysis
is incomplete.

Three features of the input language matter for what follows.

\emph{Restrictions} constrain the set of traces under consideration. A
restriction is a closed formula, and Tamarin analyzes only those traces that
satisfy it. Restrictions are how modellers express equality tests, uniqueness
of identifiers, and ordering assumptions that the rule language cannot state
directly.

\emph{Lemma annotations} change how a lemma participates in the analysis. A
lemma marked \code{[reuse]} is used as an assumption when proving the lemmas
that follow it. A lemma marked \code{[sources]} is consumed earlier still: it
refines the precomputed case distinctions that every other lemma is proved
against. A lemma marked \code{[use\_induction]} is proved by induction over the
trace rather than by the default simplification.

\emph{Wellformedness checks} run before analysis and report modelling errors:
facts used with inconsistent arity, reserved names, unbound variables. They are
the tool's own account of what it thinks is wrong with your model.

A user who writes a model, sees a clean wellformedness report, sees no lemma
reported as falsified, sees no \code{sorry} in the proof, and includes an
\code{exists-trace} lemma to confirm the model is not vacuous has done
everything the manual~\cite{TamarinManual} asks. Section~\ref{sec:last} exhibits a model that passes
every one of those checks and whose secrecy verdict is wrong.

\section{Method: one evening, and what the assistant actually did}\label{sec:method}

The audit ran on a MacBook with Tamarin 1.12.0 from Homebrew and Maude 3.5.1.
The Tamarin source tree was checked out at \code{develop} commit
\code{ef3f0468} (cabal version 1.13.0). No Haskell toolchain was installed, so
we could not build a modified Tamarin; every hypothesis had to be tested
against the stock binary using model files alone. That constraint turned out to
be useful, because it forced every finding into a form a user could reproduce.

The working loop was: read a region of source with the assistant, identify
predicates whose correctness the rest of the pipeline depends on, and for each
one ask what input would make it wrong. Then write a model that supplies that
input.

Two things are worth saying plainly about the assistant's role, since the
paper's thesis rests on it.

It was not an oracle. Several hypotheses it proposed were wrong, and the
wrongness only surfaced under test. It suggested that Tamarin's goal-ranking
oracles might prune the search unsoundly; reading \code{oracleRanking} showed
that goals the oracle omits are appended rather than dropped, so the search
cannot lose them. It suggested that a heuristic in observational-equivalence
mode might report attacks on unclosed constraint systems; the code supports
that reading, but we could not construct a witness, and \S\ref{sec:held}
reports it as an open lead rather than a finding.

What it did supply was breadth at speed. A single reader can hold one
hypothesis about one file. Working this way we examined the restriction
classifier, the wellformedness gate, the lemma-dependency machinery, the goal
manager, the injective-fact inference, the macro and predicate expanders, the
proof replay path, and the equational front end, and tested a hypothesis
against each. Most produced nothing. The ones that produced something are
\S\S\ref{sec:last}--\ref{sec:reuse}.

\subsection{Oracles that need no ground truth}\label{sec:oracles}

A recurring difficulty in auditing a verifier is that checking its answer
requires knowing the answer, which usually means trusting a second tool or a
hand proof. We used two invariants instead, both of which catch the tool
disagreeing with itself.

The first is quantifier duality. In one theory, state a property twice:
\begin{tmlisting}
lemma all_phi:            "phi"
lemma ex_notphi: exists-trace "not phi"
\end{tmlisting}
Semantics forces these to disagree. If every trace satisfies \code{phi} then no
trace satisfies its negation, and conversely. So \code{verified}/\code{verified}
and \code{falsified}/\code{falsified} are both self-contradictions, detectable
without knowing which answer is right. This invariant is what makes the defect
in \S\ref{sec:reuse} checkable by a reader who has never seen the protocol.

The second is annotation neutrality. Options documented to affect only how the
search runs, and attributes documented to be sound, must not change a verdict.
Deleting a \code{[reuse]} token, or rephrasing a restriction into a logically
equivalent form, must leave every verdict where it was. Each control in
\S\S\ref{sec:last}--\ref{sec:reuse} is an instance of this.

We built harnesses for two further invariants that did not pay off within the
evening: agreement across the sixteen documented configuration flags, and
symmetry of observational equivalence under swapping the two sides of every
\code{diff} term. Both sweeps exhausted available memory on real case studies
and were killed before completing. We ship them (\S\ref{sec:artifact}) and
claim nothing on their behalf.

\section{Defect 1: \texttt{last()} defeats the safety-formula classifier}\label{sec:last}

\subsection{The mechanism}

Tamarin's induction rule is sound only over a prefix-closed set of traces. The
implementation knows this. In \code{formulaToSystem}
(\code{lib/theory/src/Theory/Constraint/System.hs:849-856}), restrictions are
split in two, with the reason attached as a comment:

\begin{hslisting}
(safetyRestrictions, otherRestrictions) =
    partition isSafetyFormula restrictions
-- Non-safety restrictions must be added to the formula, as they
-- render the set of traces non-prefix-closed, which makes the use
-- of induction unsound.
gf2 = gconj $ gf1 : otherRestrictions
\end{hslisting}

Safety restrictions go into \code{sLemmas} and are asserted alongside the proof
obligation. Non-safety restrictions are conjoined into the obligation itself,
before induction is applied. The split is correct, and it is the fix for a
soundness bug reported in 2012 (\S\ref{sec:provenance}).

Everything therefore depends on the classifier, defined in
\code{Guarded.hs} at lines 156 to 164:

\begin{hslisting}
isSafetyFormula gf0 = null (frees [gf0]) && noExistential gf0
  where
    noExistential (GAto _)            = True
    noExistential (GGuarded Ex _ _ _) = False
    noExistential (GGuarded All _ _ gf) = noExistential gf
    noExistential (GDisj disj) = all noExistential $ getDisj disj
    noExistential (GConj conj) = all noExistential $ getConj conj
\end{hslisting}

The test is: closed, and free of existential quantifiers. The docstring above it
states the property it is standing in for, that a trace violating a safety
formula can never be extended into one satisfying it.

Closed and existential-free is not the same property. Tamarin's formula syntax
includes the atom \code{last(\#i)}, which holds when \code{\#i} is the maximal
timepoint of the trace. It is documented user-facing syntax
(\code{manual/src/011\_advanced-features.md:576}) and is parsed at
\code{Theory/Text/Parser/Formula.hs:46}. It contains no existential quantifier.
So the restriction

\begin{tmlisting}
restriction StartIsNotLast:
  "All x #i. Start(x) @ i ==> not(last(#i))"
\end{tmlisting}

passes \code{isSafetyFormula}, and Tamarin annotates it \code{// safety
formula} in its own output. It is not one. The trace $[\mathsf{Start}(x)]$
violates it, because $\mathsf{Start}$ is the final event. Its extension
$[\mathsf{Start}(x), \mathsf{Step}(x)]$ satisfies it. A violating trace extends
into a satisfying one, which is the negation of the property the docstring
names.

\subsection{The consequence}

Figure~\ref{fig:poc01} is the whole model. Both lemmas are the same formula.

\begin{figure}[t]
\begin{tmlisting}
theory LastInductionUnsound begin

rule Start: [ Fr(x) ] --[ Start(x) ]-> [ A(x) ]
rule Step:  [ A(x)  ] --[ Step(x)  ]-> [      ]

restriction StartIsNotLast:
  "All x #i. Start(x) @ i ==> not(last(#i))"

lemma NoStart_with_induction [use_induction]:
  "All x #i. Start(x) @ i ==> F"

lemma NoStart_without_induction:
  "All x #i. Start(x) @ i ==> F"

lemma StartReachable: exists-trace
  "Ex x #i. Start(x) @ i"
end
\end{tmlisting}
\caption{The minimal witness. The two \code{NoStart} lemmas are identical
formulas.}\label{fig:poc01}
\end{figure}

Tamarin's verdict:

\begin{outlisting}
NoStart_with_induction    (all-traces): verified (3 steps)
NoStart_without_induction (all-traces): falsified - found trace (2 steps)
StartReachable            (exists-trace): verified (2 steps)
/* All wellformedness checks were successful. */
\end{outlisting}

The property is false: $[\mathsf{Start}(x), \mathsf{Step}(x)]$ satisfies the
restriction and contains a $\mathsf{Start}$. Tamarin agrees when it does not use
induction. The \code{exists-trace} lemma confirms the model is not vacuous.

\subsection{The control}\label{sec:control1}

A reader might reasonably suspect the protocol, the lemma, or induction in
general. The control rules all three out. Take the same protocol and the same
lemma, and phrase the restriction as a guarded existential:

\begin{tmlisting}
restriction StartImpliesStep:
  "All x #i. Start(x) @ i ==> Ex #j. Step(x) @ j"
\end{tmlisting}

This restriction is also not prefix-closed: $[\mathsf{Start}(x)]$ violates it
and $[\mathsf{Start}(x),\mathsf{Step}(x)]$ satisfies it, exactly as before. The
difference is that it contains an existential, so \code{isSafetyFormula}
returns \code{False} and Tamarin conjoins it into the obligation. The verdict:

\begin{outlisting}
NoStart_with_induction    (all-traces): falsified - found trace (6 steps)
NoStart_without_induction (all-traces): falsified - found trace (4 steps)
\end{outlisting}

and the restriction carries no \code{// safety formula} annotation. Tamarin
handles non-prefix-closed restrictions correctly when it recognises them. The
defect is the recognition, and nothing else.

We note in passing that the most literal rendering of $\neg\mathsf{last}(i)$,
namely \tmc{Ex #j. #i < #j}, is rejected by Tamarin as an unguarded
existential, which is why the control uses an action atom as the guard.

\subsection{From a toy to a wrong security verdict}\label{sec:chain}

Figure~\ref{fig:poc01} needs an explicit \code{[use\_induction]}, which a
reviewer might notice. The annotation is not necessary. Two facts combine.

First, \code{[sources]} implies induction with no annotation from the user
(\code{lib/theory/src/ClosedTheory.hs:120}):

\begin{hslisting}
inductionHint
  | any (`elem` [SourceLemma, InvariantLemma])
        (L.get lAttributes l) = UseInduction
  | otherwise                 = AvoidInduction
\end{hslisting}

Second, \code{[sources]} lemmas are consumed during theory closing, before any
proof runs, and the case distinctions they refine are then used by every lemma
in the file (\code{CloseRule.hs:426}).

So a single misclassified restriction can wrongly prove a false sources lemma,
which then prunes legitimate cases out of the whole theory.
Figure~\ref{fig:poc02} is the result.

\begin{figure}[t]
\begin{tmlisting}
theory LastSourcesChain begin

rule Gen:  [ Fr(k) ] --[ Gen(k)  ]-> [ Out(k) ]
rule Recv: [ In(x)  ] --[ Recv(x) ]-> [        ]

restriction GenIsNotLast:
  "All k #i. Gen(k) @ i ==> not(last(#i))"

lemma bogus_sources [sources]:
  "All k #i. Gen(k) @ i ==> F"

lemma secrecy:
  "All k #i #j. Gen(k) @ i & KU(k) @ j ==> F"

lemma gen_reachable: exists-trace
  "Ex k #i. Gen(k) @ i"
end
\end{tmlisting}
\caption{The key is published by \tmc{Out(k)}. Tamarin reports its secrecy as
verified.}\label{fig:poc02}
\end{figure}

\begin{outlisting}
bogus_sources (all-traces):   verified (3 steps)
secrecy       (all-traces):   verified (2 steps)
gen_reachable (exists-trace): verified (2 steps)
/* All wellformedness checks were successful. */
\end{outlisting}

Rule \code{Gen} sends the fresh key on the public channel. The adversary learns
it immediately. Tamarin closes \code{secrecy} with \code{by solve( !KU( ~n ) @
\#j )}: the false refinement removed the only case that mattered.

Every signal the manual teaches a user to check is green. The wellformedness
report is clean. No proof contains a \code{sorry}. No lemma is reported
falsified. The \code{exists-trace} lemma confirms the model is live, so the
usual vacuity defence does not apply. Deleting only the sources lemma restores
\code{secrecy: falsified}, which confirms the mechanism.

\subsection{Provenance: the fix and the hole shipped together}\label{sec:provenance}

In August 2012 Simon Meier committed a model file titled
\code{Axioms\_and\_Induction.spthy} with the message ``example illustrating
soundness bug \#108; i.e., that our current handling of axioms and trace
induction is not sound'' (\code{a32cc3f8}). The model is still in the tree at
\code{examples/loops/Axioms\_and\_Induction.spthy}, and its comment reads
``Note that the set of traces satisfying this trace is *not* prefix-closed.''
Six days later, commit \code{9726e4c4}, ``fixed two soundness bugs'',
introduced \code{isSafetyFormula} and the partition that uses it.

We compared the function body at that commit against the current one:

\begin{outlisting}
git show 9726e4c4:src/Theory/Constraint/System/Guarded.hs
git show HEAD:lib/theory/src/Theory/Constraint/System/Guarded.hs
\end{outlisting}

The two are byte-identical. The predicate has never been modified in fourteen
years.

The question this raises is whether \code{last} was reachable from surface
syntax in 2012, or whether the hole opened later when the atom was exposed. It
was reachable. The same commit that introduced the fix contains, at
\code{src/Theory/Text/Parser.hs:340}:

\begin{hslisting}
[ Last <$> try (symbol "last" *> parens nodevarTerm) <?> "last atom"
\end{hslisting}

The fix for bug \#108 and the syntax that defeats it were in the same commit,
on the same day. The 2012 counterexample used an existential, which the new
predicate caught. The \code{last} phrasing of the same idea was already
expressible and was not caught, and has not been caught since.

\subsection{Scope, stated against our own interest}\label{sec:scope}

The defect is real and it is reachable from documented syntax. It does not
invalidate any published Tamarin result, and we think a paper that implied
otherwise would deserve to be dismissed.

We scanned every model shipped with Tamarin: 1042 files under
\code{examples/}. Thirty-six contain the string \code{last(}. In every one, the
occurrence is inside Tamarin's own machine-generated proof text, saved into an
analyzed output, of the form

\begin{outlisting}
solve( (last(#i)) | (Ex #j. Update(k,r) @ #j & !(last(#j)) & #j < #i) )
\end{outlisting}

None is a user-written restriction. Zero models use \code{last()} in the
position that triggers the bug. The scanner is
\code{attacks/tools/lastscan.py}.

So the accurate claim is that this is a live unsoundness in documented syntax
that currently affects no published result. It is a trap that has been armed
for fourteen years and that nobody has stepped on, which is a different and
less comfortable statement than either ``Tamarin's results are wrong'' or
``this does not matter''.

\subsection{The fix}

One line, in \code{Guarded.hs}:

\begin{hslisting}
noExistential (GAto (Last _)) = False
\end{hslisting}

\code{toInductionHypothesis} already refuses formulas containing \code{last}
(\code{Guarded.hs:619}, error string ``formula not last-free''), so
reclassifying makes induction inapplicable to theories carrying such a
restriction, which is the correct behaviour.

\section{Defect 2: the vacuity checks have not run since PR \#725}\label{sec:vacuity}

The most common modelling error in Tamarin is a lemma that quantifies over an
action fact no rule emits. Such a lemma is vacuously true and Tamarin reports
it \code{verified}. Practitioners are taught to guard against this by hand, but
Tamarin also ships two checks for it, \code{inexistentActions} and
\code{inexistentActionsRestrictions}
(\code{Theory/Tools/Wellformedness.hs:706-734}), fully implemented and with
error messages written.

They do not run. Here is the call site
(\code{Wellformedness.hs:578-583}):

\begin{hslisting}
-- | Report on facts usage. Skip checks on non-existant actions
-- if `incompleteMSRs` is True.
factReports incompleteMSRs thy =
    concat [ reservedReport, ..., factUsage, factLhsOccurNoRhs ]
    ++ concat [ inexistentActions ++ inexistentActionsRestrictions
              | incompleteMSRs ]
\end{hslisting}

The list comprehension \emph{includes} the checks when \code{incompleteMSRs} is
\code{True}, which is the inverse of what its own docstring says. And the sole
production call site fixes the flag (\code{src/Main/TheoryLoader.hs:602}):

\begin{hslisting}
incompleteMSRs = False
  -- TODO how do we know if we do not have all MSRs due to translation?
\end{hslisting}

Either reading condemns the code. Under the docstring's intent the guard is
backwards; under the guard's literal meaning the flag is set wrongly. In both
cases the checks are unreachable in the shipped tool.

\subsection{Witness}

\begin{figure}[t]
\begin{tmlisting}
theory VacuousLemmaTypo begin

rule Setup: [ Fr(k)  ] --[ SecretKey(k) ]-> [ Key(k) ]
rule Leak:  [ Key(k) ] --[ Leaked(k)    ]-> [ Out(k) ]

// The rules emit SecretKey. This lemma says SecrettKey.
lemma secrecy_typo:
  "All k #i. SecrettKey(k) @ i ==> not(Ex #j. K(k) @ j)"

lemma secrecy_intended:
  "All k #i. SecretKey(k) @ i ==> not(Ex #j. K(k) @ j)"
end
\end{tmlisting}
\caption{A one-character typo. Tamarin reports the vacuous lemma as verified
and the wellformedness report as clean.}\label{fig:poc03}
\end{figure}

\begin{outlisting}
secrecy_typo     (all-traces): verified (2 steps)
secrecy_intended (all-traces): falsified - found trace (3 steps)
/* All wellformedness checks were successful. */
\end{outlisting}

Running with \code{--quit-on-warning} changes nothing, because no warning is
produced: the check that would produce it never executes.

\subsection{Provenance}

\code{git log -L 583,583} on the file attributes the guard to commit
\code{98c3887f}, ``Extend ProVerif export module (\#725)''. We diffed the
commit directly. Before it:

\begin{hslisting}
factReports :: OpenTranslatedTheory -> WfErrorReport
factReports thy = concat
    [ reservedReport, reservedFactNameRules, freshFactArguments
    , specialFactsUsage, factUsage, factLhsOccurNoRhs
    , inexistentActions, inexistentActionsRestrictions ]
\end{hslisting}

After it, the last two members moved behind the guard shown above. A refactor
of an unrelated export backend silently disabled the tool's only automatic
guard against its most common modelling error, and no test caught the
regression.

\section{Defect 3: unproven assumptions are admitted as facts}\label{sec:reuse}

\subsection{The mechanism}

\code{gatherReusableLemmas} (\code{lib/theory/src/CloseRule.hs:181-188}) decides
which lemmas are admitted as assumptions:

\begin{hslisting}
gatherReusableLemmas kind = do
    LemmaItem lem <- previousItems
    guard $    lemmaSourceKind lem <= kind
            && ReuseLemma `elem` L.get lAttributes lem
            && AllTraces == L.get lTraceQuantifier lem
            && L.get lName lem `notElem` L.get pcHiddenLemmas ctxt
            && "ALL" `notElem` L.get pcHiddenLemmas ctxt
    return $ formulaToGuarded_ $ L.get lFormula lem
\end{hslisting}

The filter checks the attribute, the trace quantifier, the source kind and the
hide-list. It does not check whether the lemma was proved. We looked for a
proof-status gate elsewhere on the path and there is none: the
\code{CompleteProof} constructor is consumed only in \code{src/Web/Theory.hs},
where it colours a cell in the interactive interface.

The comment beside the insertion states the justification:

\begin{hslisting}
-- Note that it is OK to add reusable lemmas directly to the system,
-- as they do not change the considered set of traces. This is the key
-- difference between lemmas and restrictions.
\end{hslisting}

That is sound under the precondition that the lemma holds. Nothing enforces the
precondition.

\subsection{Tamarin contradicting itself}

The admitted formulas go into \code{sLemmas}, which constrains the constraint
system for both trace quantifiers. A false assumption therefore does two things
at once: it closes all-traces branches that should stay open, and it prunes the
witness search for exists-trace lemmas. Figure~\ref{fig:poc08} exploits both.

\begin{figure}[t]
\begin{tmlisting}
theory ReusePoisonsBothDirections begin

rule Gen: [ Fr(k) ] --[ Secret(k) ]-> [ Out(k) ]

lemma helper [reuse]:
  "All k #i #j. Secret(k) @ i & KU(k) @ j ==> F"

lemma target:
  "All k #i. Secret(k) @ i ==> not(Ex #j. KU(k) @ j)"

lemma live: exists-trace
  "Ex k #i #j. Secret(k) @ i & KU(k) @ j"
end
\end{tmlisting}
\caption{\code{helper} and \code{live} are exact negations of one
another.}\label{fig:poc08}
\end{figure}

\begin{outlisting}
helper (all-traces):   falsified - found trace (3 steps)
target (all-traces):   verified (2 steps)
live   (exists-trace): falsified - no trace found (2 steps)
\end{outlisting}

Read the first and third lines together. \code{helper} is
\tmc{All k #i #j. Secret(k)@i & KU(k)@j ==> F} and \code{live} is
\tmc{Ex k #i #j. Secret(k)@i & KU(k)@j}. The second is the negation of the
first. ``Falsified'' on \code{helper} asserts that such a trace exists.
``Falsified'' on \code{live} asserts that no such trace exists. Tamarin affirms
and denies the same proposition in one summary block.

This is the defect we consider most serious, and not because its consequences
are the worst. It is the one that requires nothing of the reader. No knowledge
of the protocol, no judgement about what the right answer is, and no trust in
us. Two lines of the tool's own output are inconsistent with each other.

Why is the inconsistency visible at all, given that both poisoned verdicts
might have been consistent with one another? Because \code{previousItems}
draws only from lemmas that precede the current one, so \code{helper} never
reuses itself and receives the correct verdict, while every later lemma is
evaluated under the refuted assumption. The pair \code{target}/\code{live}, by
contrast, comes back \code{verified}/\code{falsified}, which is mutually
consistent and therefore invisible to the duality oracle of
\S\ref{sec:oracles}. Detecting this class requires pairing against the reuse
lemma itself.

\subsection{The control}

Deleting the single token \code{[reuse]} from Figure~\ref{fig:poc08}, changing
nothing else, gives

\begin{outlisting}
helper (all-traces):   falsified - found trace (3 steps)
target (all-traces):   falsified - found trace (3 steps)
live   (exists-trace): verified (3 steps)
\end{outlisting}

all three correct. This fixes the ground truth, so the contradiction above
resolves as follows: \code{live} is true, \code{target} is false, and Tamarin's
verdicts for both are wrong.

\subsection{The same defect on a hotter path}

\code{[sources]} lemmas are admitted the same way and consumed earlier, during
theory closing, with no failure path if their proof later fails. The blast
radius is the whole theory rather than one lemma:

\begin{outlisting}
bogus_sources (all-traces): falsified - found trace (2 steps)
target        (all-traces): verified (2 steps)
\end{outlisting}

The auto-generated sources lemmas produced by \code{--auto-sources} take the
same unguarded path. We checked whether this is realised in practice: all
sixteen auto-sources case studies in \code{case-studies-regression/} prove
their generated \code{AUTO\_typing} lemma, so the risk there is structural
rather than current.

\subsection{On invalidation}

Tamarin has a proof status for exactly this hazard.
\code{Theory/Proof.hs:406} defines \code{InvalidatedProof}, documented as ``The
proof has been Invalidated (eg. by editing a reuse lemma)''. The constructor is
set in exactly one place in the codebase: \code{src/Web/Handler.hs:255}, the
interactive web interface.

We want to state the consequence carefully, because the obvious framing is
wrong. The manual documents this invalidation in its interactive-mode section
(\code{manual/src/007\_property-specification.md:229-234}), so batch mode never
promised it. The correct statement is not that a safety net fails to fire. It
is that the safety net exists only in the graphical interface, while batch
mode, which is what continuous integration and every published artifact use,
has no equivalent and no proof-status gate anywhere on the reuse path.

\section{What we could not break}\label{sec:held}

This section is longer than any of the preceding three, which is the point.

We report it for two reasons. Negative results locate a defect: knowing that
the solver held while the classifier failed is what turns five bug reports into
a claim about where assurance effort is misallocated. And a paper that reported
only hits would deserve the suspicion that it went looking for a conclusion.

\paragraph{The constraint solver and term algebra.} The duality oracle of
\S\ref{sec:oracles} was run over fifteen models covering exclusive-or
(self-cancellation, one-time pad, chained pads, the \code{zero} constant),
multisets, Diffie-Hellman, subterms, natural numbers, hashing, and symmetric
encryption. No contradictions. Controls in the same battery return opposite
verdicts as they should, so the oracle discriminates.

\paragraph{Injective facts and stateful reasoning.} \code{simpInjectiveFactEqMon}
(\code{Simplify.hs:534-660}) is the newest and most aggressive inference in the
solver: it derives equalities and ordering constraints from inferred monotonic
behaviour of fact arguments, which is the dangerous direction, since it adds
constraints rather than removing them. It held on constant positions, strictly
increasing positions, natural-number counters, non-fresh identifiers, and
SAPIC locked and unlocked state cells.

\paragraph{Stored-proof replay.} Tamarin serialises proofs into the model file,
and we expected replay to be a soft target. It is not. Three separate forgeries
against \code{case-studies-regression/Tutorial\_analyzed.spthy}, namely
weakening a lemma statement while keeping its proof, replacing a proof with a
bogus one-step \code{by contradiction}, and relabelling a case, were all caught
and downgraded to \code{analysis incomplete}. Published proof artifacts do get
re-validated on load.

\paragraph{Goal-ranking oracles.} \code{oracleRanking}
(\code{ProofMethod.hs:604-621}) returns \code{ranked ++ remaining}. Goals the
external oracle omits are appended rather than dropped, so a hostile or buggy
oracle costs time and not soundness. The only termination path,
\code{quitOnEmpty}, emits a \code{Sorry} and is reported as an incomplete proof.

\paragraph{Predicate expansion.} \code{expandFormula}
(\code{Theory/Syntactic/Predicate.hs:94-105}) shifts de Bruijn indices
correctly. A predicate body binding \code{y}, used at a call site that also
binds \code{y}, gives results identical to a non-colliding control. Macro
expansion, by contrast, is not hygienic: a macro body variable is captured by a
same-named rule variable, silently, with a clean wellformedness report. We do
not report that as a defect, because the genuinely dangerous case, a macro
injecting an unbound variable, is caught by \code{unboundReport}, and because
we could not turn the capture into a wrong verdict.

\paragraph{Goal management.} The \code{openGoals} filter drops \code{KU} goals
it treats as always constructible (\code{Goals.hs:70-84}). We checked whether
the dropped goals are genuinely handled elsewhere; they are, decomposed into
sub-goals by \code{insertAction} (\code{Reduction.hs:293-350}).

\paragraph{Equational front end.} Non-terminating equations are rejected at
parse time. Equations that are not subterm-convergent produce an explicit
warning naming the offending equation. Precomputation limits (\code{-s},
\code{-c}) leave sources under-refined when hit, which costs completeness and
not soundness.

\paragraph{An open lead.} In observational-equivalence mode, \code{DiffAttack}
(\code{ProofMethod.hs:373-387}) can report an attack on a constraint system it
never closed:

\begin{hslisting}
guard (solved || (trivial sys' && notContradictory))
\end{hslisting}

The load-bearing assumption is that a goal over pairwise-distinct message
variables is always satisfiable. It is asserted in a comment, and it has failed
before: commit \code{37c7285a} is titled ``fixed bug when mirroring XOR leading
to false attacks'', and \code{07021243} narrowed the same check with the
message ``otherwise unsound''. We could not construct a witness. Our two
attempts, a reflexivity check on \code{diff(t,t)} and a symmetry sweep over
real case studies, produced one completed passing case and one sweep that
exhausted memory before finishing. This lead is open, not cleared. It is also
the most promising direction for the one result this paper lacks, a false
attack rather than a missed one.

\subsection{The shape of the result set}

Every confirmed defect makes Tamarin too permissive. Not one makes it too
strict. For a tool whose output is used to certify deployed protocols this is
the worse direction, and it is consistent with the pattern: each defect is a
place where a check was skipped, never a place where an inference was drawn
incorrectly. The inferences are, as far as we could determine in an evening,
correct.

\section{Disclosure}\label{sec:disclosure}

We attempted coordinated disclosure and could not find a channel to use.

At the time of writing, the Tamarin repository has no security policy. There is
no \code{SECURITY.md} at the repository root or under \code{.github/}, so
GitHub's advisory interface at
\code{github.com/tamarin-prover/tamarin-prover/security} offers no reporting
path. \code{CONTRIBUTING.md} documents pull requests and issue etiquette and
does not mention vulnerabilities, security reports, or private disclosure. The
\code{.github/} directory contains two workflow files and nothing else.

For most projects the absence of a policy is unremarkable. For a tool whose
output is cited in support of the security of deployed protocols, a defect that
turns \code{falsified} into \code{verified} is closer to a vulnerability than
to an ordinary bug, and the absence of a route to report one privately is worth
naming.

We are therefore publishing directly, with complete reproductions, and we would
encourage the maintainers to adopt a policy. Our first recommendation in
\S\ref{sec:lessons} is that one be added, because it costs a file and it is the
precondition for everything else in that list.

\section{What we think this means}\label{sec:lessons}

\subsection{Sophistication and assurance are not the same axis}

Tamarin's constraint-solving procedure is a serious piece of work with a
published metatheory, and our evening produced no evidence against it. The
defects are all in code that any competent programmer could have audited: a
five-line predicate, an inverted list comprehension, a filter missing a
conjunct. This is not a coincidence, and we do not think it is specific to
Tamarin.

Attention follows interest. The constraint solver is where the research is, so
it gets the papers, the proofs, and the scrutiny. The restriction classifier is
plumbing. But the classifier is what decides whether the solver's soundness
theorem applies to your model, and a soundness theorem about a set of traces
you are not analyzing is not a soundness theorem about your protocol. The
guardrails carry the theorem across to the user, and they had not been audited
to anything like the depth of the thing they carry.

\subsection{A regression test would have caught two of three}

The 2012 counterexample for bug \#108 is still in the tree. Nobody added its
\code{last}-phrased sibling. The vacuity checks were unconditional before PR
\#725 and gated afterwards, and no test noticed, because the test suite
compares end-to-end verdicts on case studies and no case study depends on those
checks firing.

Both are cheap to prevent. A property test asserting that
\code{isSafetyFormula} implies prefix-closure on generated formulas would have
caught the first at any point in fourteen years. A single model with a
deliberately misspelled action fact, asserting that the wellformedness report
is non-empty, would have caught the second at the moment of the refactor.

\subsection{On the cost of an audit}

The claim in our title is meant literally, and we want to be careful with it in
both directions.

An evening was enough. It was enough because reading unfamiliar Haskell with an
assistant that can follow a call graph across twenty files is much faster than
reading it alone, and because the hypotheses that mattered were cheap to test:
a ten-line model file and a sub-second run. The assistant was wrong repeatedly,
and every finding here survived because it was reduced to a reproduction rather
than to an argument. That is the part of the method we would defend most
strongly, and it is not new; it is just newly cheap.

An evening was also not enough for plenty. We did not complete the
configuration-agreement sweep. We did not complete the observational-equivalence
symmetry sweep. We did not seriously examine the SAPIC translation, the export
backends, or the accountability machinery. We have one open lead and no false
attack. A determined effort over a week, or the same method applied by people
who know the codebase, would find more.

The lesson we would draw is not about Tamarin, and it is certainly not that
tools like it should be trusted less in favour of nothing. It is that the
expected cost of auditing them has dropped by something like an order of
magnitude, and the field has not yet adjusted its habits to that fact. Every
verifier in common use, ours included, deserves an evening of this. We expect
most of them will yield something, and we expect that to be good for all of us.

\subsection{Recommendations}

\begin{enumerate}
  \item Add a \code{SECURITY.md} with a private reporting address.
  \item Fix \code{isSafetyFormula} by rejecting \code{Last} atoms, one line, and
    add the \code{last}-phrased sibling of the 2012 counterexample to the
    regression suite.
  \item Repair the \code{incompleteMSRs} guard so the vacuity checks execute,
    and add a model whose wellformedness report is asserted non-empty.
  \item Refuse to admit a \code{[reuse]} or \code{[sources]} lemma whose proof
    is not closed, or propagate a taint marker into dependent lemmas so the
    summary reads \code{verified (assuming: helper)} rather than
    \code{verified}.
  \item Add property tests for the predicates the soundness argument depends
    on, separately from the end-to-end case-study suite, since the latter
    cannot see a check that stopped running.
\end{enumerate}

\section{Artifact}\label{sec:artifact}

The artifact is seven Tamarin model files, a runner, and two analysis scripts.
Five models reproduce defects and two are controls. \code{run.sh} executes all
seven, checks each verdict against the expected pattern, and exits non-zero if
any defect stops reproducing, so it doubles as a regression test against a
patched build. Every model completes in under a second.

\code{tools/lastscan.py} produces the corpus scan of \S\ref{sec:scope}.
\code{tools/consistency\_oracle.py} implements the duality oracle of
\S\ref{sec:oracles} and contains the fifteen-case battery of \S\ref{sec:held}.
We also ship the two harnesses whose sweeps did not complete
(\code{differ.py}, \code{diffsym.py}); they are included because the invariants
they encode are the right ones, and we claim no result on their behalf.

All findings were verified against \code{tamarin-prover 1.12.0} with Maude
3.5.1, and every implicated code path was compared across \code{1.12.0..HEAD}
and is unchanged, so the defects are present on \code{develop} at
\code{ef3f0468}.

\section{Related work}\label{sec:related}

Soundness defects in verification tools have been reported before, and the
genre is well established: implementations of decision procedures have been
found to disagree with their own metatheory in SMT
solvers~\cite{SMTFuzz,SMTDelta} and in C compilers~\cite{Csmith}, generally by
differential or metamorphic testing rather than by proof. Our duality oracle
(\S\ref{sec:oracles}) is a metamorphic relation in that tradition, specialised
to a tool that answers questions about trace sets.

Within protocol verification, comparisons between tools have surfaced
discrepancies that turned out to be tool defects rather than modelling
differences, and Tamarin's own history records several: the two soundness bugs
fixed in \code{9726e4c4}, the mirroring bug in \code{37c7285a}, and the
narrowing in \code{07021243}, each found by the developers themselves. What we
add is not a new class of defect but a data point about cost: the same kind of
finding, obtained without domain expertise in the tool, in one evening.

\section{Conclusion}

We found five defects in Tamarin in one evening. Three are serious: an
unsound induction rule reachable from documented syntax and open since 2012, a
pair of wellformedness checks that have not executed since an unrelated
refactor, and a lemma-dependency mechanism that admits refuted assumptions and
can be made to affirm and deny the same proposition in a single run.

We also spent that evening failing to break the constraint solver, the term
algebra, the monotonicity inference, the goal manager, and proof replay. That
failure is the more interesting half of the result. The hard parts of Tamarin
are in good shape. The checks that decide whether those hard parts apply to
your model are not, and nobody had looked, because looking used to be
expensive and now it is not.

None of this makes Tamarin a bad tool. We would still reach for it. It does
suggest that the confidence the field places in tools of this class has been
calibrated against their theory rather than against their implementations, and
that the gap between the two is now cheap enough to measure that there is no
longer an excuse for leaving it unmeasured. That applies to every tool in this
space, including our own. It seemed worth an evening to find out. We would
recommend the exercise.


\begin{thebibliography}{99}

\bibitem{TamarinCAV}
S.~Meier, B.~Schmidt, C.~Cremers, and D.~Basin.
\newblock The TAMARIN Prover for the Symbolic Analysis of Security Protocols.
\newblock In \emph{Computer Aided Verification (CAV)}, LNCS 8044, pp. 696--701,
  2013.
\newblock \url{https://doi.org/10.1007/978-3-642-39799-8_48}.

\bibitem{TamarinCSF}
B.~Schmidt, S.~Meier, C.~Cremers, and D.~Basin.
\newblock Automated Analysis of Diffie-Hellman Protocols and Advanced Security
  Properties.
\newblock In \emph{IEEE Computer Security Foundations Symposium (CSF)},
  pp. 78--94, 2012.
\newblock \url{https://doi.org/10.1109/CSF.2012.25}.

\bibitem{TamarinManual}
The Tamarin Team.
\newblock \emph{Tamarin-Prover Manual: Security Protocol Analysis in the
  Symbolic Model}.
\newblock \url{https://tamarin-prover.com/manual/}. Accessed August 2026.

\bibitem{DolevYao}
D.~Dolev and A.~C.~Yao.
\newblock On the Security of Public Key Protocols.
\newblock \emph{IEEE Transactions on Information Theory}, 29(2):198--208, 1983.
\newblock \url{https://doi.org/10.1109/TIT.1983.1056650}.

\bibitem{ProVerif}
B.~Blanchet.
\newblock An Efficient Cryptographic Protocol Verifier Based on Prolog Rules.
\newblock In \emph{IEEE Computer Security Foundations Workshop (CSFW)},
  pp. 82--96, 2001.
\newblock \url{https://doi.org/10.1109/CSFW.2001.930138}.

\bibitem{TLS13}
C.~Cremers, M.~Horvat, J.~Hoyland, S.~Scott, and T.~van~der~Merwe.
\newblock A Comprehensive Symbolic Analysis of TLS 1.3.
\newblock In \emph{ACM Conference on Computer and Communications Security
  (CCS)}, pp. 1773--1788, 2017.
\newblock \url{https://doi.org/10.1145/3133956.3134063}.

\bibitem{FiveG}
D.~Basin, J.~Dreier, L.~Hirschi, S.~Radomirovic, R.~Sasse, and V.~Stettler.
\newblock A Formal Analysis of 5G Authentication.
\newblock In \emph{ACM Conference on Computer and Communications Security
  (CCS)}, pp. 1383--1396, 2018.
\newblock \url{https://doi.org/10.1145/3243734.3243846}.

\bibitem{EMV}
D.~Basin, R.~Sasse, and J.~Toro-Pozo.
\newblock The EMV Standard: Break, Fix, Verify.
\newblock In \emph{IEEE Symposium on Security and Privacy (S\&P)},
  pp. 1766--1781, 2021.
\newblock \url{https://doi.org/10.1109/SP40001.2021.00037}.

\bibitem{Csmith}
X.~Yang, Y.~Chen, E.~Eide, and J.~Regehr.
\newblock Finding and Understanding Bugs in C Compilers.
\newblock In \emph{ACM SIGPLAN Conference on Programming Language Design and
  Implementation (PLDI)}, pp. 283--294, 2011.
\newblock \url{https://doi.org/10.1145/1993498.1993532}.

\bibitem{SMTFuzz}
D.~Winterer, C.~Zhang, and Z.~Su.
\newblock On the Unusual Effectiveness of Type-Aware Operator Mutations for
  Testing SMT Solvers.
\newblock \emph{Proceedings of the ACM on Programming Languages}, 4(OOPSLA),
  Article 193, pp. 193:1--193:25, 2020.
\newblock \url{https://doi.org/10.1145/3428261}.

\bibitem{SMTDelta}
R.~Brummayer and A.~Biere.
\newblock Fuzzing and Delta-Debugging SMT Solvers.
\newblock In \emph{7th International Workshop on Satisfiability Modulo Theories
  (SMT)}, pp. 1--5, 2009.
\newblock \url{https://doi.org/10.1145/1670412.1670413}.

\end{thebibliography}

\end{document}
